\documentclass{beamer}
\usepackage{datetime}

\usepackage[french]{babel}
\usepackage[T1]{fontenc}

\usepackage{amsmath, amssymb, amsfonts, amsthm}
\usepackage{booktabs}

\usetheme{metropolis}           % Use metropolis theme

\title{Théorie des modèles}
\date{04/06/26}
\author{Lucas Gaebele \& Adan Herlaut}
\institute{Université Côte d'Azur}
\begin{document}
  \maketitle
  \section{Introduction express à la théorie des modèles}
\begin{frame}{Prélude}
    La Théorie des modèles étudie les relations entre les énoncés formels (syntaxe) et les structures mathématiques qui les réalisent (sémantique).
\end{frame}

\begin{frame}{Langages}
    Un langage $\mathcal{L}$ est une collection de symboles divisée en trois types : les fonctions, les relations et les constantes.
    \begin{example}
    Écrire les axiomes de la théorie des anneaux suppose un langage :
    \begin{enumerate}
    \item Des fonctions : $+, \times$
    \item Des constantes : $0, 1$
    \item Aucune relation
    \end{enumerate}
    On note $\mathcal{L}_{\text{ring}} = \{+, \times, 0, 1\}$ le langage des anneaux.
    \end{example}
\end{frame}

\begin{frame}{Formules}
On se donne un langage $\mathcal{L}$. A l'aide des symboles donnés par la Logique ($\wedge, \implies, =, \forall$...) et ceux donnés par $\mathcal{L}$,
nous pouvons écrire des $\mathcal{L}$-formules.
\begin{example}
Reprenons $\mathcal{L}_{\text{ring}} = \{+, \times, 0, 1\}$.
On peut alors écrire les formules : 
\begin{enumerate}
\item $\phi_1 : \forall x,y,z [ x \times (y+z) = x \times y + x \times z ]$
\item $ \phi_2 : \forall x [ \neg(x = 0) \implies \exists y, x \times y = 1 ]$
\end{enumerate}
\end{example}
\end{frame}

\begin{frame}{Théories et modèles}
\begin{definition}
Soit $\mathcal{L}$ un langage. Une $\mathcal{L}$-théorie $\mathcal{T}$ est un ensemble de $\mathcal{L}$-formules.
Un modèle $\mathcal{M}$ de $\mathcal{T}$ est une structure \textbf{mathématique} où toutes les formules de $\mathcal{T}$ sont \textbf{vraies} dans $\mathcal{M}$.
\end{definition}
\begin{example}
Muni de $\mathcal{L}_{\text{ring}} = \{+, \times, 0, 1\}$, on peut écrire la \textbf{théorie des anneaux} $\mathcal{T}$. $\phi_1$ est un élément de $\mathcal{T}$.
$\mathbb{Z}$ est un modèle de $\mathcal{T}$.
\end{example}
\end{frame}

\begin{frame}{Théories et modèles}
\begin{definition}
\begin{enumerate}
\item Soit $\mathcal{T}$ une $\mathcal{L}$-théorie. $\mathcal{T}$ est \textbf{satisfiable} s'il existe un modèle $\mathcal{M}$ de $\mathcal{T}$.
\item Soit $\phi$ une $\mathcal{L}$-formule, on note $\mathcal{T} \models \phi$ si $\phi$ est vraie dans tous les modèles de $\mathcal{T}$.
\end{enumerate}
\end{definition}
\begin{example}
Soit $\mathcal{T}$ la théorie des anneaux. On a $\mathcal{T} \models \phi_1$, mais $\mathcal{T} \not\models \phi_2$.
\end{example}
\end{frame}

\begin{frame}{Théorème de compacité}
\begin{theorem}
Soit $T$ une $\mathcal{L}$-théorie. $T$ est satisfiable si et seulement si tout sous-ensemble fini de $T$ est satisfiable.
\end{theorem}
\end{frame}

\begin{frame}{Théorie des corps algébriquement clos}
Soit $\mathcal{L}_{\text{ring}} = \{+, \times, 0, 1\}$. Formalisons la \textbf{théorie des corps algébriquement clos}.
On rajoute à la théorie des corps les $\mathcal{L}_{ring}$-formules suivantes :
\begin{align*}
 \phi_n : \forall a_0, \ldots, a_{n-1}, \exists x [x^n+a_{n-1}\cdot x^{n-1} + \cdots + a_1 \cdot x + a_0 = 0].
\end{align*}
pour tout $n \in \mathbb{N^*}$. Ceci donne la théorie des corps algébriquement clos, notée $ACF$.
\end{frame}

\begin{frame}{Théorie des corps algébriquement clos}
Soit $p \in \mathbb{P}$ et la $\mathcal{L}_{\text{ring}}$-formule $\chi_p : \underbrace{1 + \cdots 1 = 0}_\text{p}$.\\On définit :
\begin{enumerate}
\item $ACF_p = ACF \cup \chi_p$ la théorie des corps algébriquement clos de caractéristique $p$.
\item $ACF_0 = ACF \cup \{\neg\chi_p \mid p \in \mathbb{P}\}$ la théorie des corps algébriquement clos de caractéristique nulle.
\end{enumerate}
\end{frame}

  \section{De la complétude au théorème d'Ax-Grothendieck}
 % Def complétude, ACF complete
\begin{frame}{Complétude}
    \begin{definition}[Théorie complète]
        Une $\mathcal{L}$-théorie $T$ est dite \textbf{complète} si pour toute $\mathcal{L}$-phrase $\phi$, on a soit :
        \[ T \models \phi \quad \text{ou} \quad T \models \neg\phi \]
    \end{definition}
    
    \pause
    \begin{theorem}[Complétude de $ACF_p$]
        Soit $p \in \{0\} \cup \mathcal{P}$ la caractéristique. La théorie $ACF_p$ est \textbf{complète}.
    \end{theorem}
\end{frame}


% Principe de Lefschetz
\begin{frame}{Le Principe de Lefschetz}
    La complétude et la \textbf{compacité} permettent de transférer les énoncés entre les caractéristiques:
    
    \begin{table}
        \centering
        \begin{tabular}{ll}
            \toprule
            \textbf{Énoncé $\phi$ vrai sur...} & \textbf{Équivaut à...} \\
            \midrule
            Le corps $\mathbb{C}$ & Vrai sur tout $K \models ACF_0$ \\
            Tout $K \models ACF_0$ & Vrai sur $ACF_p$ pour $p \gg 0$ \\
            \bottomrule
        \end{tabular}
    \end{table}

    \pause
    \begin{exampleblock}{Idée de la preuve}
        Si $ACF_0 \models \phi$, alors il existe un sous-ensemble fini d'axiomes $\Sigma \subset ACF_0$ tel que $\Sigma \models \phi$. Comme $\Sigma$ est fini, il n'exclut qu'un nombre fini de caractéristiques $p$.
    \end{exampleblock}
\end{frame}

% Énoncé Ax-Grothendieck
\begin{frame}{Le Théorème d'Ax-Grothendieck}
    \begin{theorem}[Ax-Grothendieck, 1966]
        Soit $f : \mathbb{C}^n \to \mathbb{C}^n$ une application polynomiale. 
        Si $f$ est \textbf{injective}, alors elle est \textbf{surjective}.
    \end{theorem}

    \pause
    
    \begin{exampleblock}{Stratégie de la preuve}
        Utiliser le principe de Lefschetz pour se ramener à un problème simple sur les corps finis.
    \end{exampleblock}
\end{frame}

% Preuve AG 1
\begin{frame}{Preuve : Étape 1 - Traduction de l'énoncé}
    On fixe $n$ (variables) et $d$ (degré maximum). L'application polynomiale
    $f$ est déterminée par ses coefficients $\bar{c}$. \\
    % exemple au tableau
    \pause
    On peut écrire une phrase $\phi_{n,d}$ dans le langage $\mathcal{L}_{ring}$:
    \[ \phi_{n,d} : \forall \bar{c} \, (\text{Inj}(\bar{c}) \to \text{Surj}(\bar{c})) \]
    
    \pause
    Par principe de Lefschetz,pour prouver $\phi_{n,d}$ sur $\mathbb{C}$ il
    suffit de le prouver sur $\overline{\mathbb{F}_p}$ pour tout $p$.

\end{frame}

% Preuve AG 2
\begin{frame}{Preuve : Étape 2 - Raisonnement par l'absurde}
    On travaille sur $\overline{\mathbb{F}_p}$. Soit $f$ injective. Supposons par l'absurde qu'il existe $y \notin \text{Im}(f)$.
    
    \pause
    \begin{itemize}
        \item Soit $K$ le sous-corps de $\overline{\mathbb{F}_p}$ engendré par les coefficients de $f$ et les composantes de $y$.
        \pause
        \item $K$ est une extension finie de $\mathbb{F}_p$, donc \textbf{$K$ est un corps fini}.
    \end{itemize}
    
    \pause
    \begin{alertblock}{Contradiction}
        $f$ induit une application injective $f|_K : K^n \to K^n$. \\

        \pause

        Or, pour un ensemble fini, \textbf{injection $\iff$ surjection}. \\

        Donc $y$ doit être dans l'image de $f$.
    \end{alertblock}
\end{frame}

  
  \section{De l'élimination des quantificateurs au Nullstellensatz}

    
  % Def elimination des quantificateurs, th Tarski

  \begin{frame}
\begin{definition}[Élimination des quantificateurs]
        Une théorie $T$ admet l'élimination des quantificateurs si pour toute formule $\varphi(\bar{x})$, il existe une formule $\theta(\bar{x})$ \textbf{sans quantificateurs} telle que :
        \[ T \models \forall \bar{x} \, (\varphi(\bar{x}) \leftrightarrow \theta(\bar{x})) \]
    \end{definition}
  \end{frame}

  \pause 
  \begin{theorem}[Tarski, 1951]
        La théorie \textbf{ACF} admet l'élimination des quantificateurs.
    \end{theorem}

    % geo alg
  \begin{frame}{Le Dictionnaire Algèbre-Géométrie}
    Pour lier la l'algèbre à la géométrie, on définit deux correspondances fondamentales entre les points de $k^n$ et les polynômes de $K[\overline{x}]$.

    \pause
    \begin{block}{1. De l'Algèbre vers la Géométrie: l'ensemble des zéros communs}
        Soit $S \subset k[\bar{x}]$ un ensemble de polynômes. L'ensemble de leurs zéros communs est:
        \[ V(S) = \{ x \in k^n \mid \forall f \in S, f(x) = 0 \} \]
        
    \end{block}
    \pause
    \begin{block}{2. De la Géométrie vers l'Algèbre: l'idéal annulateur}
        Soit $X \subset k^n$ un ensemble de points. L'ensemble des polynômes s'annulant sur $X$ est:
        \[ I(X) = \{ f \in k[\bar{x}] \mid \forall x \in X, f(x) = 0 \} \]
    \end{block}
\end{frame}

% 
\begin{frame}{Une correspondance parfaite?}
    On vérifie facilement que $V(I(X)) = X$. \\
    \vspace{0.5cm}
    \textbf{Question :} A-t-on $I(V(I)) = I$ pour tout idéal $I$?

    \pause
    \begin{alertblock}{Contre-exemple}
        Plaçons-nous dans $\mathbb{C}[X]$ et considérons l'idéal $I = (X^2)$.
        \pause
        \begin{enumerate}
            \item L'ensemble des zéros est: $V(I) = \{ 0 \}$. \pause
            \item L'idéal annulateur de ce point est: $I(\{0\}) = (X)$. \pause
            \item Conclusion: $(X) \neq (X^2)$.
        \end{enumerate}
    \end{alertblock}
\end{frame}

\begin{frame}{Nullstellensatz}
    \begin{definition}[Radical d'un idéal]
        Le radical de $I$, noté $\sqrt{I}$, est l'ensemble des éléments dont une puissance appartient à $I$:
        \[ \sqrt{I} = \{ f \in k[\bar{x}] \mid \exists m \ge 1, f^m \in I \} \]
    \end{definition}

    \pause
    \begin{theorem}[Nullstellensatz de Hilbert]
        Soit $k$ un corps algébriquement clos et $I$ un idéal de $K[\bar{x}]$. Alors:
        \[ I(V(I)) = \sqrt{I} \]
    \end{theorem}
    \pause
    \textit{Vérification sur notre exemple :} $\sqrt{(X^2)} = (X)$.
\end{frame}

\begin{frame}{Esquisse de la preuve}
    L'inclusion $\sqrt{I} \subseteq I(V(I))$ est immédiate. Prouvons l'autre sens.
    
    \pause
    \vspace{0.3cm}
    Soit $f \in I(V(I))$. On introduit une nouvelle variable $Y$ et l'idéal :
    \[ J = (I, 1 - Yf) \subset K[\overline{x}, Y] \]
    
    \pause
    \begin{enumerate}
        \item \textbf{Absence de zéros:} Si $x \in V(I)$, alors $f(x) = 0$, donc $1 - Y(0) = 1 \neq 0$. Ainsi, $V(J) = \emptyset$.
        
        \pause
        \item \textbf{Nullstellensatz faible:} Par l'élimination des quantificateurs, $J$ n'est pas un idéal propre, donc $1 \in J$.
        
        \pause
        \item \textbf{Identité de Bézout:} Il existe $A_i, B$ tels que 
        \[ 1 = \sum P_i A_i + (1-Yf)B \]
        \pause
        \item \textbf{Évaluation:} En évaluant en $Y = 1/f(\overline{x})$ et en chassant les dénominateurs, on obtient $f^m \in I$, d'où $f \in \sqrt{I}$.
    \end{enumerate}
\end{frame}

\end{document}